\subsubsection{Task, benchmark and metrics}
\label{sec:task}

\paragraph{Task.} The input is a generated clip $V$, a target string $s$ and a hint box $b$ that marks where the text is. In our benchmark the target is the quoted sign text of the generation prompt, and the hint box is produced automatically by OCR localisation and tracking of the generated text; the same box is given to every method, and \steer uses it as its editing mask. The output is a clip $\hat V$ of the same size and length in which the text inside $b$ reads $s$ and nothing else has changed. Unlike scene text editing, the source text is usually not a valid word---it is a near miss produced by the generator---and in about half of the clips it is already correct, so ``do nothing'' is often the right answer. We call the task \emph{repair} to stress both aspects.

\paragraph{Benchmark.} We generated short clips of street scenes, shop fronts and screens with the text-to-video mode of VACE-14B~\cite{jiang2025vace}, from 120 prompts that ask for a specific sign text on one of seven carrier materials, with ten seeds each (832$\times$464 pixels, 49 frames at 15\,fps). Of 1{,}178 usable clips, an OCR screen confirmed the requested text in only 24\%. The screen reads every fourth frame with EasyOCR and keeps, per frame, the best normalised edit similarity between a recognised string and the target; a clip counts as correct when at least 80\% of these frames reach 0.8, and as garbled when it does not but its mean similarity is at least 0.3, so that it contains a recognisable attempt at the target. This rule selects the 526 garbled clips and drops the other 371, which show no such attempt. The benchmark therefore measures the repair of a near miss, not reconstruction without any source text, which we study separately (\cref{sec:exp_neg}). The screen is stricter than our evaluation reader, which judges \rOrigRead\% of them correct, so the benchmark contains both repair and do-nothing cases. The 526 clips come from 117 prompts with up to ten seeds each, 40 target strings (20 English, 20 Chinese, from one to four words or characters, \eg \str{BUS STOP}, \str{SALE 50\%}, \zh{小心地滑}), and seven carrier materials: neon (93 clips), signboard (86), paper (96), glass (76), screen (65), metal (56) and wood (54); 295 clips are English and 231 Chinese. Every clip has a hint box per frame. In the evaluation job of \cref{tab:main}, the reader (below) judges \rOrigCorrect{} originals (\rOrigRead\%) correct and \rGarbled{} garbled. Forty clips were looked at while developing the selective-editing components; they are listed in the supplement and excluded in a separate report.

\subsubsection{Metrics}
\label{sec:metrics}
\paragraph{Readability.} A vision-language reader, Qwen2.5-VL-7B~\cite{bai2025qwen25vl}, transcribes the text inside the hint box (15\% padding) on 5 frames spread over the clip. A clip counts as \emph{correct} if at least 4 of the 5 transcriptions match the target exactly after removing whitespace and line breaks, lower-casing, and folding traditional to simplified Chinese; readability is the share of correct clips, and NED similarity the mean normalised edit similarity over the 5 frames. The reader is independent of \method, whose gates use a different OCR system. All outputs that are compared are read in the same job and software environment, and a second reader from a different model family re-reads the main comparison with the same protocol (supplement).\footnote{We learned this the hard way: one version of the reader produced garbled transcriptions for about 8.5\% of the inputs on one cluster and none on another; comparing outputs read in different environments inflated an early result (supplement).}

\paragraph{Repairs and breakages.} Readability alone hides two opposite effects, so we also count \emph{fixed} clips (garbled $\to$ correct), \emph{broken} clips (correct $\to$ wrong) and their difference, the \emph{net} repairs.

\paragraph{Locality.} \emph{BG PSNR} is computed between output and input on all pixels outside the hint box dilated by half its height, over 49 frames. \emph{Collateral change} is the mean absolute grey-level change, on the reference frame, inside the text neighbourhood (box grown by half its height) but outside both the old glyphs and the new glyphs (each dilated by 0.15 box heights); it isolates damage that is not part of writing the new text, and it requires the method's glyph mask, so we report it for \method only.

\paragraph{Appearance.} The original glyphs are wrong but their material is the reference. On three frames, we estimate an ink mask inside the box for every video and compare, in CIELAB, the ink colour, the ink--ground contrast, the halo around the strokes (normalised by contrast), the edge softness and the shading inside the strokes with those of the original. Each feature is divided by its median over all method--clip pairs of the evaluation run and the five are averaged; lower is closer. The normalisation makes values comparable only within one run, which is why every table compares variants from a single run. The ink segmentation is not seeded, so re-scoring an identical video changes its distance by \rRetestSD{} (standard deviation per clip); we interpret only differences well above this noise (supplement).

\paragraph{Temporal stability.} The carrier's frame-to-frame motion is estimated once per clip on the original (SIFT, RANSAC homography) and applied to every output, so camera motion cancels. \emph{Flicker} is the mean absolute motion-compensated grey-level change inside the box, \emph{jitter} the same for gradient magnitudes (stroke position), and \emph{boundary jump} the change across the frame pairs at which a 17-frame chunked model switches chunks, relative to all other pairs (1 means no seam; used for the learned refiner in \cref{sec:exp_refiner}).

\paragraph{Human judgements.} Blind pairwise studies show the original next to two versions in random left/right order, with the target known to the raters; questions are asked separately (\eg material fidelity, artifacts, flicker), and ties are allowed. Study designs, sampling and rater quality control are in the supplement.
